\section{Proving equivalence of X25519 in C and Coq}
\label{C-Coq}

In this section we first describe the global structure of our proof.
Then we discuss techniques used to prove the equivalence between the
Clight description of TweetNaCl and Coq functions producing similar behaviors.

In order to prove the correctness of X25519 in TweetNaCl code, we use VST to prove
that the code matches our functional specification of \texttt{crypto\_scalarmult}
(abreviated as CSM) in Coq. Then, we prove that our specification of the scalar
multiplication matches the mathematical definition of elliptic curves and Theorem 2.1 by
Bernstein~\cite{Ber06}. Figure~\ref{tk:ProofOverview} shows a graph of dependencies
of the proofs considered. The mathematical proof of our specification is presented
in Section~\ref{sec3-maths}.

\begin{figure}[h]
  \centering
  \begin{tikzpicture}[textstyle/.style={black, anchor= south west, align=center, minimum height=0.45cm, text centered, font=\scriptsize}]


  % adapt line thickness and line width, if needed

   \path [thick, dashed] (2.5,1) edge +(0,-6.75);
   \draw (2.5,1) node[anchor=north east] {\sref{sec:Coq-RFC}};
   \draw (2.5,1) node[anchor=north west] {\sref{sec:C-Coq}};
   \path [thick, dashed] (0,-5.75) edge +(8.5,0);
   \draw (8.5,-5.75) node[anchor=north east] {\sref{sec:maths}};

  % SECTION III

  % Definition of RFC
  \begin{scope}[yshift=-3 cm,xshift=0 cm]
   \draw (0,0) -- (0.4, 0.4) -- (1.75,0.4) -- (1.75,0) -- cycle;
   \draw (0,0) -- (1.75,0) -- (1.75,-1) -- (0, -1) -- cycle;
   \draw (0.3,-0.05) node[textstyle] {Definition};
   \draw (0.875,-0.5) node[textstyle, anchor=center] {\texttt{RFC}};
   \draw (1.75,-1) node[textstyle, anchor = south east] {\color{doc@lstfunctions}{.V}};
  \end{scope}

   % SECTION IV

  % C code
  \begin{scope}[yshift=-0.25 cm,xshift=2.75 cm]
    \draw (0,0) -- (0.4, 0.4) -- (1.25,0.4) -- (1.25,0) -- cycle;
    \draw (0,0) -- (1.25,0) -- (1.25,-1.25) -- (0, -1.25) -- cycle;
    \draw (0.3,-0.05) node[textstyle] {code};
    \draw (0.675,-0.5) node[textstyle, anchor=center] {\texttt{Prog}};
    \draw (1.25,-1.25) node[textstyle, anchor = south east] {\color{doc@lstfunctions}{.C}};
  \end{scope}

  % V code
  \begin{scope}[yshift=-0.25 cm,xshift=5 cm]
    \draw (0,0) -- (0.4, 0.4) -- (1.25,0.4) -- (1.25,0) -- cycle;
    \draw (0,0) -- (1.25,0) -- (1.25,-1.25) -- (0, -1.25) -- cycle;
    \draw (0.3,-0.05) node[textstyle] {code};
    \draw (0.675,-0.5) node[textstyle, anchor=center] {\texttt{Prog}};
    \draw (1.25,-1.25) node[textstyle, anchor = south east] {\color{doc@lstfunctions}{.V}};
    % \draw (1.25,0)  node[anchor=south east, inner sep=0pt] {\includegraphics[width=.0125\textwidth]{img/coq_logo.png}};
  \end{scope}

  \path [thick, dotted, double, ->] (4,-0.75) edge (5, -0.75);

  % VST Spec
  \begin{scope}[yshift=-3 cm,xshift=2.75 cm]
    \draw (0,0) -- (0.4, 0.4) -- (2.75,0.4) -- (2.75,0) -- cycle;
    \draw (0,0) -- (2.75,0) -- (2.75,-2) -- (0, -2) -- cycle;
    \draw (0.3,-0.05) node[textstyle] {Specification};
    \draw (1.375,-1) node[textstyle, anchor=center, align=left] {
      {\color{doc@lstnumbers}\textbf{Pre}}:\\
      ~~\texttt{s}$[\!\!\{$\texttt{n}$\}\!\!]\leftarrow n \in \N$,\\
      % ~~$n \in$ \TNaCles{u8[32]},\\
      ~~\texttt{s}$[\!\!\{$\texttt{p}$\}\!\!]\leftarrow P \in E(\F{p^2})$\\
      % ~~$P \in$ \TNaCles{u8[32]}\\
      {\color{doc@lstdirective}\textbf{Post}}:\\
      ~~\texttt{s}$[\!\!\{$\texttt{q}$\}\!\!]\leftarrow$~\texttt{RFC}$(n,P)$};
      \draw (2.75,-2) node[textstyle, anchor = south east] {\color{doc@lstfunctions}{.V}};
  \end{scope}

  % VST Theorem
  \begin{scope}[yshift=-3 cm,xshift=6 cm]
    \draw (0,0) -- (0.4, 0.4) -- (2.5,0.4) -- (2.5,0) -- cycle;
    \draw (0,0) -- (2.5,0) -- (2.5,-1.25) -- (0, -1.25) -- cycle;
    \draw (0.3,-0.05) node[textstyle] {Proof};
    \draw (1.25,-0.5) node[textstyle, anchor=center] {\{{\color{doc@lstnumbers}\textbf{Pre}}\} \texttt{Prog} \{{\color{doc@lstdirective}\textbf{Post}}\}};
    \draw (2.5,0) node[textstyle, anchor = south east] {\color{doc@lstfunctions}{\checkmark}};
    \draw (2.5,-1.25) node[textstyle, anchor = south east] {\color{doc@lstfunctions}{.V}};
  \end{scope}

  \path [thick, double] (5.625,-1.5) edge [out=-90, in=90] (5.625, -2.5);
  \path [thick, double, ->] (5.625, -2.5) edge [out=-90, in=180] (6, -3.5);
  \path [thick, double, ->] (5.5,-3.75) edge [out=0, in=180] (6, -3.75);



  % SECTION V

  % Spec of Curve nP
  \begin{scope}[yshift=-7.5 cm,xshift=0 cm]
    \draw (0,0) -- (0.4, 0.4) -- (1.75,0.4) -- (1.75,0) -- cycle;
    \draw (0,0) -- (1.75,0) -- (1.75,-1) -- (0, -1) -- cycle;
    \draw (0.3,-0.05) node[textstyle] {Definition};
    \draw (0.875,-0.5) node[textstyle, anchor=center] {$n \cdot P$};
    \draw (1.75,-1) node[textstyle, anchor = south east] {\color{doc@lstfunctions}{.V}};
  \end{scope}

  % Correctness Theorem
  \begin{scope}[yshift=-7 cm,xshift=6 cm]
    \draw (0,0) -- (0.4, 0.4) -- (2.5,0.4) -- (2.5,0) -- cycle;
    \draw (0,0) -- (2.5,0) -- (2.5,-1.5) -- (0, -1.5) -- cycle;
    \draw (0.3,-0.05) node[textstyle] {Proof};
    \draw (1.25,-0.75) node[textstyle, anchor=center] {{\color{doc@lstnumbers}\textbf{Pre}} $\implies$\\$\text{\texttt{RFC}}(n,P) = n \cdot P$};
    \draw (2.5,0) node[textstyle, anchor = south east] {\color{doc@lstfunctions}{\checkmark}};
    \draw (2.5,-1.5) node[textstyle, anchor = south east] {\color{doc@lstfunctions}{.V}};
  \end{scope}

  \path [thick, double, ->] (1.75,-3.5) edge  [out=0, in=-180] (2.75, -3.5);
  \path [thick, double] (1.75,-3.5) edge [out=0, in=90] (2.25, -4);
  \path [thick, double] (2.25, -4) edge [out=-90, in=90] (2.25, -7);
  \path [thick, double] (2.25, -7) edge [out=-90, in=-180] (3, -7.75);
  \path [thick, double, ->] (3, -7.75) edge [out=0, in=-180] (6, -7.75);
  \path [thick, double, ->] (1.75,-8) edge [out=0, in=-180] (6, -8);


\end{tikzpicture}

  \caption{Overview of the proof}
  \label{tk:ProofOverview}
\end{figure}


Verifying \texttt{crypto\_scalarmult} also implies to verify all the functions
subsequently called: \texttt{unpack25519}; \texttt{A}; \texttt{Z}; \texttt{M};
\texttt{S}; \texttt{car25519}; \texttt{inv25519}; \texttt{set25519}; \texttt{sel25519};
\texttt{pack25519}.

We prove the implementation of Curve25519 is \textbf{sound}, \ie:
\begin{itemize}
\item absence of access out-of-bounds of arrays (memory safety).
\item absence of overflows/underflow on the arithmetic.
\end{itemize}
We also prove that TweetNaCl's code is \textbf{correct}:
\begin{itemize}
\item Curve25519 is correctly implemented (we get what we expect).
\item Operations on \texttt{gf} (\texttt{A}, \texttt{Z}, \texttt{M}, \texttt{S})
are equivalent to operations ($+,-,\times,x^2$) in \Zfield.
\item The Montgomery ladder computes a scalar multiplication between a natural
number and a point.
\end{itemize}

In order to prove the soundness and correctness of \texttt{crypto\_scalarmult},
we first create a skeleton of the Montgomery ladder with abstract operations which
can be instantiated over lists, integers, field elements...
A high level specification (over a generic field $\K$) allows us to prove the
correctness of the ladder with respect to the curves theory.
This high-level specification does not rely on the parameters of Curve25519.
By instantiating $\K$ with $\Zfield$, and the parameters of Curve25519 ($a = 486662, b = 1$),
we define a mid-level specification.
Additionally we also provide a low-level specification close to the \texttt{C} code
(over lists of $\Z$). We show this specification to be equivalent to the
\textit{semantic version} of C (\texttt{CLight}) with VST.
This low level specification gives us the soundness assurance.
By showing that operations over instances ($\K = \Zfield$, $\Z$, list of $\Z$) are
equivalent, we bridge the gap between the low level and the high level specification
with Curve25519 parameters.
As such, we prove all specifications to equivalent (Fig.\ref{tk:ProofStructure}).
This garantees us the correctness of the implementation.

\begin{figure}[h]
  \begin{tikzpicture}[textstyle/.style={black, anchor= south west, align=center}]

    \filldraw[draw=orange!10!doc@lstbackground, fill=doc@lstbackground, thick] (0.25,0.5) rectangle (4.5,5.5);
    % node[textstyle, anchor=west, draw=yellow, fill=yellow!20, thick, minimum width=5.5cm,minimum height=5cm] {};

    \draw (4.5,5.5)  node[anchor=north east, inner sep=0pt] (russell) {\includegraphics[width=.03\textwidth]{img/coq_logo.png}};

    \draw (0.5,-1) node [textstyle, anchor=west, draw=black, thick, minimum width=3cm,minimum height=0.5cm] (longlong) {\texttt{long long[16]}};
    \draw (0,-1) node [anchor=east] (longlongdef) {\texttt{C} code};
    \draw (2.25,-0.1) node [anchor=west] (app) {\texttt{clightgen}};


    \draw (0.5,1) node [textstyle, anchor=west, draw=black, thick, minimum width=3cm,minimum height=0.5cm] (clight) {\texttt{tptr tlong}};
    \draw (0,1) node [anchor=east] (clightdef) {Clight};

    \draw [thick, ->] (longlong.north) -- (clight.south);

    \begin{scope}[yshift=1 cm,xshift=0 cm]

      \draw (0.5,2) node [textstyle, anchor=west, draw=black, thick, minimum width=3cm,minimum height=0.5cm] (ll) {\texttt{list} $\Z$};
      \draw (0,2) node [anchor=east] (shn) {Low Level};

      \draw [thick,double, <->, >=implies] (clight.north) -- (ll.south);
      \draw (2.25,1) node[anchor=west, inner sep=0pt] (chain) {\includegraphics[width=.07\textwidth]{img/chain.png}};


      \draw (0.5,3) node [textstyle, anchor=west, draw=black, thick, minimum width=3cm,minimum height=0.5cm] (ml) {$\Zfield$};
      \draw (0,3) node [anchor=east] (shn) {Mid Level};

      \draw[thick,double, <->, >=implies] (ll.north) -- (ml.south);

      \draw (0.5,4) node [textstyle, anchor=west, draw=black, thick, minimum width=3cm,minimum height=0.5cm] (hl) {$\K$};
      \draw (0,4) node [anchor=east] (shn) {High Level};

      \draw[thick,double, <-, >=implies] (ml.north) -- (hl.south);
    \end{scope}
\end{tikzpicture}

  \caption{Structural construction of the proof}
  \label{tk:ProofStructure}
\end{figure}



\subsection{Correctness Specification}

We show the soundness of TweetNaCl by proving the following specification matches a pure Coq function.
This defines the equivalence between the Clight representation and a Coq definition of the ladder.

\begin{CoqVST}
Definition crypto_scalarmult_spec :=
DECLARE _crypto_scalarmult_curve25519_tweet
WITH
  v_q: val, v_n: val, v_p: val, c121665:val,
  sh : share,
  q : list val, n : list Z, p : list Z
(*------------------------------------------*)
PRE [ _q OF (tptr tuchar),
     _n OF (tptr tuchar),
     _p OF (tptr tuchar) ]
PROP (writable_share sh;
      Forall (fun x => 0 <= x < 2^8) p;
      Forall (fun x => 0 <= x < 2^8) n;
      Zlength q = 32; Zlength n = 32;
      Zlength p = 32)
LOCAL(temp _q v_q; temp _n v_n; temp _p v_p;
      gvar __121665 c121665)
SEP  (sh [{ v_q }] <<(uch32)-- q;
      sh [{ v_n }] <<(uch32)-- mVI n;
      sh [{ v_p }] <<(uch32)-- mVI p;
      Ews [{ c121665 }] <<(lg16)-- mVI64 c_121665)
(*------------------------------------------*)
POST [ tint ]
PROP (Forall (fun x => 0 <= x < 2^8) (CSM n p);
      Zlength (CSM n p) = 32)
LOCAL(temp ret_temp (Vint Int.zero))
SEP  (sh [{ v_q }] <<(uch32)-- mVI (CSM n p);
      sh [{ v_n }] <<(uch32)-- mVI n;
      sh [{ v_p }] <<(uch32)-- mVI p;
      Ews [{ c121665 }] <<(lg16)-- mVI64 c_121665
\end{CoqVST}

In this specification we state as preconditions:
\begin{itemize}
  \item[] \VSTe{PRE}: \VSTe{_p OF (tptr tuchar)}\\
  The function \texttt{crypto\_scalarmult} takes as input three pointers to
  arrays of unsigned bytes (\VSTe{tptr tuchar}) \VSTe{_p}, \VSTe{_q} and \VSTe{_n}.
  \item[] \VSTe{LOCAL}: \VSTe{temp _p v_p}\\
  Each pointer represent an address \VSTe{v_p},
  \VSTe{v_q} and \VSTe{v_n}.
  \item[] \VSTe{SEP}: \VSTe{sh [{ v_p $\!\!\}\!\!]\!\!\!$ <<(uch32)-- mVI p}\\
  In the memory share \texttt{sh}, the address \VSTe{v_p} points
  to a list of integer values \VSTe{mVI p}.
  \item[] \VSTe{PROP}: \VSTe{Forall (fun x => 0 <= x < 2^8) p}\\
  In order to consider all the possible inputs, we assumed each
  elements of the list \texttt{p} to be bounded by $0$ included and $2^8$
  excluded.
  \item[] \VSTe{PROP}: \VSTe{Zlength p = 32}\\
  We also assumed that the length of the list \texttt{p} is 32. This defines the
  complete representation of \TNaCle{u8[32]}.
\end{itemize}

As Post-condition we have:
\begin{itemize}
  \item[] \VSTe{POST}: \VSTe{tint}\\
  The function \texttt{crypto\_scalarmult} returns an integer.
  \item[] \VSTe{LOCAL}: \VSTe{temp ret_temp (Vint Int.zero)}\\
  The returned integer has value $0$.
  \item[] \VSTe{SEP}: \VSTe{sh [{ v_q $\!\!\}\!\!]\!\!\!$ <<(uch32)-- mVI (CSM n p)}\\
  In the memory share \texttt{sh}, the address \VSTe{v_q} points
  to a list of integer values \VSTe{mVI (CSM n p)} where \VSTe{CSM n p} is the
  result of the \VSTe{crypto_scalarmult} over \VSTe{n} and \VSTe{p}.
  \item[] \VSTe{PROP}: \VSTe{Forall (fun x => 0 <= x < 2^8) (CSM n p)}\\
  \VSTe{PROP}: \VSTe{Zlength (CSM n p) = 32}\\
  We show that the computation for \VSTe{CSM} fits in  \TNaCle{u8[32]}.
\end{itemize}

This specification shows that \VSTe{crypto_scalarmult} in C computes the same
result as \VSTe{CSM} in Coq provided that inputs are within their respective
bounds.
By converting those array of 32 bytes into their respective little-endian value
we prove the correctness of \VSTe{crypto_scalarmult} (Theorem \ref{CSM-correct})
in Coq (for the sake of simplicity we do not display the conversion in the theorem).
\begin{theorem}
\label{CSM-correct}
For all $n \in \N, n < 2^{255}$ and where the bits 1, 2, 5 248, 249, 250
are cleared and bit 6 is set, for all $P \in E(\F{p^2})$,
for all $p \in \F{p}$ such that $P.x = p$,
there exists $Q \in E(\F{p^2})$ such that $Q = \cdot P$ where $Q.x = q$ and $q$ = \VSTe{CSM} $n$ $p$.
\end{theorem}
A more complete description in Coq of Theorem \ref{CSM-correct} with the associated conversions
is as follow:
\begin{lstlisting}[language=Coq]
Theorem Crypto_Scalarmult_Correct:
  forall (n p:list Z) (P:mc curve25519_Fp2_mcuType),
  Zlength n = 32 ->
  Zlength p = 32 ->
  Forall (fun x => 0 <= x /\ x < 2^8) n ->
  Forall (fun x => 0 <= x /\ x < 2^8) p ->
  Fp2_x (ZUnpack25519 (ZofList 8 p)) = P#x0 ->
  ZofList 8 (Crypto_Scalarmult n p) =
    (P *+ (Z.to_nat (Zclamp (ZofList 8 n)))) _x0.
\end{lstlisting}

\subsection{Number Representation and C Implementation}

As described in Section \ref{preliminaries:B}, numbers in \TNaCle{gf} are represented
in base $2^{16}$ and we use a direct mapping to represent that array as a list
integers in Coq. However in order to show the correctness of the basic operations,
we need to convert this number as a full integer.
\begin{definition}
Let \Coqe{ZofList} : $\Z \rightarrow \texttt{list} \Z \rightarrow \Z$, a parametrized map by $n$ betwen a list $l$ and its
it's little endian representation with a base $2^n$.
\end{definition}
We define it in Coq as:
\begin{lstlisting}[language=Coq]
Fixpoint ZofList {n:Z} (a:list Z) : Z :=
  match a with
  | [] => 0
  | h :: q => h + (pow 2 n) * ZofList q
  end.
\end{lstlisting}
We define a notation where $n$ is $16$.
\begin{lstlisting}[language=Coq]
Notation "Z16.lst A" := (ZofList 16 A).
\end{lstlisting}
We also define a notation to do the modulo, projecting any numbers in $\Zfield$.
\begin{lstlisting}[language=Coq]
Notation "A :GF" := (A mod (2^255-19)).
\end{lstlisting}
% Remark that this representation is different from \Coqe{Zmodp}.
However the equivalence between operations over $\Zfield$ and $\F{p}$ is easily proven.

Using these two definitions, we proved intermediates lemmas such as the correctness of the
multiplication \Coqe{M} where \Coqe{M} replicate the computations and steps done in C.
\begin{lemma}
For all list of integers \texttt{a} and \texttt{b} of length 16 representing
$A$ and $B$ in $\Zfield$, the number represented in $\Zfield$ by the list \Coqe{(M a b)}
is equal to $A \times B \bmod \p$.
\end{lemma}
And seen in Coq as follows:
\begin{Coq}
Lemma mult_GF_Zlength :
  forall (a:list Z) (b:list Z),
  Zlength a = 16 ->
  Zlength b = 16 ->
   (Z16.lst (M a b)) :GF =
   (Z16.lst a * Z16.lst b) :GF.
\end{Coq}

\subsection{Inversions in \Zfield}

We define a Coq version of the inversion mimicking
the behavior of \TNaCle{inv25519} over \Coqe{list Z}.
\begin{lstlisting}[language=Ctweetnacl]
sv inv25519(gf o,const gf a)
{
  gf c;
  int i;
  set25519(c,a);
  for(i=253;i>=0;i--) {
    S(c,c);
    if(i!=2 && i!=4) M(c,c,a);
  }
  FOR(i,16) o[i]=c[i];
}
\end{lstlisting}
We specify this with 2 functions: a recursive \Coqe{pow_fn_rev} to to simulate the for loop and a simple
\Coqe{step_pow} containing the body. Note the off by one for the loop.
\begin{lstlisting}[language=Coq]
Definition step_pow (a:Z) (c g:list Z) : list Z :=
  let c := Sq c in
    if a <>? 1 && a <>? 3
    then M c g
    else c.

Function pow_fn_rev (a:Z) (b:Z) (c g: list Z)
  {measure Z.to_nat a} : (list Z) :=
  if a <=? 0
    then c
    else
      let prev := pow_fn_rev (a - 1) b c g in
        step_pow (b - 1 - a) prev g.
\end{lstlisting}

This \Coqe{Function} requires a proof of termination. It is done by proving the
well-foundness of the decreasing argument: \Coqe{measure Z.to_nat a}. Calling
\Coqe{pow_fn_rev} 254 times allows us to reproduce the same behavior as the \texttt{Clight} definition.
\begin{lstlisting}[language=Coq]
Definition Inv25519 (x:list Z) : list Z :=
  pow_fn_rev 254 254 x x.
\end{lstlisting}
Similarily we define the same function over $\Z$.
\begin{lstlisting}[language=Coq]
Definition step_pow_Z (a:Z) (c:Z) (g:Z) : Z :=
  let c := c * c in
  if a <>? 1 && a <>? 3
    then c * g
    else c.

Function pow_fn_rev_Z (a:Z) (b:Z) (c:Z) (g: Z)
  {measure Z.to_nat a} : Z :=
  if (a <=? 0)
    then c
    else
      let prev := pow_fn_rev_Z (a - 1) b c g in
        step_pow_Z (b - 1 - a) prev g.

Definition Inv25519_Z (x:Z) : Z :=
  pow_fn_rev_Z 254 254 x x.
\end{lstlisting}
And prove their equivalence in $\Zfield$.
\begin{lstlisting}[language=Coq]
Lemma Inv25519_Z_GF : forall (g:list Z),
  length g = 16 ->
  (Z16.lst (Inv25519 g)) :GF =
  (Inv25519_Z (Z16.lst g)) :GF.
\end{lstlisting}
In TweetNaCl, \TNaCle{inv25519} computes an inverse in $\Zfield$.
It uses the Fermat's little theorem by doing an exponentiation to $2^{255}-21$.
This is done by applying a square-and-multiply algorithm. The binary representation
of $p-2$ implies to always do multiplications aside for bits 2 and 4.
To prove the correctness of the result we can use multiple strategies such as:
\begin{itemize}
  \item Proving it is a special case of square-and-multiply algorithm applied to $2^{255}-21$.
  \item Unrolling the for loop step-by-step and applying the equalities
  $x^a \times x^b = x^{(a+b)}$ and $(x^a)^2 = x^{(2 \times a)}$. We prove that
  the resulting exponent is $2^{255}-21$.
\end{itemize}
We use the second method for the benefits of simplicity. However it requires to
apply the unrolling and exponentiation formulas 255 times. This can be automated
in Coq with tacticals such as \Coqe{repeat}, but it generates a proof object which
will take a long time to verify.

\subsection{Speeding up with Reflections}

In order to speed up the verification, we use a technique called reflection.
It provides us with flexibility, for example, we don't need to know the number of
times nor the order in which the lemmas needs to be applied (chapter 15 in \cite{CpdtJFR}).

The idea is to \textit{reflect} the goal into a decidable environment.
We show that for a property $P$, we can define a decidable boolean property
$P_{bool}$ such that if $P_{bool}$ is \Coqe{true} then $P$ holds.
$$reify\_P : P_{bool} = true \implies P$$
By applying $reify\_P$ on $P$ our goal become $P_{bool} = true$.
We then compute the result of $P_{bool}$. If the decision goes well we are
left with the tautology $true = true$.

To prove that the \Coqe{Inv25519_Z} is computing $x^{2^{255}-21}$,
we define a Domain Specific Language.
\begin{definition}
Let \Coqe{expr_inv} denote an expression which is either a term;
a multiplication of expressions; a squaring of an expression or a power of an expression.
And Let \Coqe{formula_inv} denote an equality between two expressions.
\end{definition}
\begin{lstlisting}[language=Coq]
Inductive expr_inv :=
  | R_inv : expr_inv
  | M_inv : expr_inv -> expr_inv -> expr_inv
  | S_inv : expr_inv -> expr_inv
  | P_inv : expr_inv -> positive -> expr_inv.

Inductive formula_inv :=
  | Eq_inv : expr_inv -> expr_inv -> formula_inv.
\end{lstlisting}
The denote functions are defined as follows:
\begin{lstlisting}[language=Coq]
Fixpoint e_inv_denote (m:expr_inv) : Z :=
  match m with
  | R_inv     =>
    term_denote
  | M_inv x y =>
    (e_inv_denote x) * (e_inv_denote y)
  | S_inv x =>
    (e_inv_denote x) * (e_inv_denote x)
  | P_inv x p =>
    pow (e_inv_denote x) (Z.pos p)
  end.

Definition f_inv_denote (t : formula_inv) : Prop :=
  match t with
  | Eq_inv x y => e_inv_denote x = e_inv_denote y
  end.
\end{lstlisting}
All denote functions also take as an argument the environment containing the variables.
We do not show it here for the sake of readability.
Given that an environment, \Coqe{term_denote} returns the appropriate variable.
With such Domain Specific Language we have the equality between:
\begin{lstlisting}[backgroundcolor=\color{white}]
f_inv_denote
 (Eq_inv (M_inv R_inv (S_inv R_inv))
         (P_inv R_inv 3))
  = (x * x^2 = x^3)
\end{lstlisting}
On the right side, \Coqe{(x * x^2 = x^3)} depends on $x$. On the left side,
\texttt{(Eq\_inv (M\_inv R\_inv (S\_inv R\_inv)) (P\_inv R\_inv 3))} does not depend on $x$.
This allows us to use computations in our decision precedure.

We define \Coqe{step_inv} and \Coqe{pow_inv} to mirror the behavior of
\Coqe{step_pow_Z} and respectively \Coqe{pow_fn_rev_Z} over our DSL and
we prove their equality.
\begin{lstlisting}[language=Coq]
Lemma step_inv_step_pow_eq :
  forall (a:Z) (c:expr_inv) (g:expr_inv),
  e_inv_denote (step_inv a c g) =
  step_pow_Z a (e_inv_denote c) (e_inv_denote g).

Lemma pow_inv_pow_fn_rev_eq :
  forall (a:Z) (b:Z) (c:expr_inv) (g:expr_inv),
  e_inv_denote (pow_inv a b c g) =
  pow_fn_rev_Z a b (e_inv_denote c) (e_inv_denote g).
\end{lstlisting}
We then derive the following lemma.
\begin{lemma}
\label{reify}
With an appropriate choice of variables,
\Coqe{pow_inv} denotes \Coqe{Inv25519_Z}.
\end{lemma}

In order to prove formulas in \Coqe{formula_inv},
we have the following a decidable procedure.
We define \Coqe{pow_expr_inv}, a function which returns the power of an expression.
We then compare the two values and decide over their equality.
\begin{Coq}
Fixpoint pow_expr_inv (t:expr_inv) : Z :=
  match t with
  | R_inv   => 1
  (* power of a term is 1. *)
  | M_inv x y =>
    (pow_expr_inv x) + (pow_expr_inv y)
  (* power of a multiplication is
     the sum of the exponents. *)
  | S_inv x =>
    2 * (pow_expr_inv x)
  (* power of a squaring is the double
     of the exponent. *)
  | P_inv x p =>
    (Z.pos p) * (pow_expr_inv x)
  (* power of a power is the multiplication
     of the exponents. *)
  end.

Definition decide_e_inv (l1 l2:expr_inv) : bool :=
  (pow_expr_inv l1) ==? (pow_expr_inv l2).

Definition decide_f_inv (f:formula_inv) : bool :=
  match f with
  | Eq_inv x y => decide_e_inv x y
  end.
\end{Coq}
We prove our decision procedure correct.
\begin{lemma}
\label{decide}
For all formulas $f$, if the decision over $f$ returns \Coqe{true},
then the denoted equality by $f$ is true.
\end{lemma}
Which is formalized as:
\begin{Coq}
Lemma decide_formula_inv_impl :
  forall (f:formula_inv),
  decide_f_inv f = true ->
  f_inv_denote f.
\end{Coq}
By reification to over DSL (lemma \ref{reify}) and by applying our decision (lemma \ref{decide}).
we proved the following theorem.
\begin{theorem}
\Coqe{Inv25519_Z} computes an inverse in \Zfield.
\end{theorem}
\begin{Coq}
Theorem Inv25519_Z_correct :
  forall (x:Z),
  Inv25519_Z x = pow x (2^255-21).
\end{Coq}

From \Coqe{Inv25519_Z_GF} and \Coqe{Inv25519_Z_correct}, we conclude the
functionnal correctness of the inversion over \Zfield.
\begin{corollary}
\Coqe{Inv25519} computes an inverse in \Zfield.
\end{corollary}
\begin{Coq}
Corollary Inv25519_Zpow_GF :
  forall (g:list Z),
  length g = 16 ->
  Z16.lst (Inv25519 g) :GF  =
  (pow (Z16.lst g) (2^255-21)) :GF.
\end{Coq}

\subsection{Packing and other Applications of Reflection}

We prove the functional correctness of \Coqe{Inv25519} with reflections.
This technique can also be used where proofs requires some computing or a small and
finite domain of variable to test e.g. for all $i$ such that $0 \le i < 16$.
Using reflection we prove that we can split the for loop in \TNaCle{pack25519} into two parts.
\begin{lstlisting}[language=Ctweetnacl]
for(i=1;i<15;i++) {
  m[i]=t[i]-0xffff-((m[i-1]>>16)&1);
  m[i-1]&=0xffff;
}
\end{lstlisting}
The first loop is computing the substraction while the second is applying the carrying.
\begin{lstlisting}[language=Ctweetnacl]
for(i=1;i<15;i++) {
  m[i]=t[i]-0xffff
}
for(i=1;i<15;i++) {
  m[i]=m[i]-((m[i-1]>>16)&1);
  m[i-1]&=0xffff;
}
\end{lstlisting}
This loop separation allows simpler proofs. The first loop is seen as the substraction of a number in \Zfield.
We then prove that with the iteration of the second loop, the number represented in $\Zfield$ stays the same.
This leads to the proof that \TNaCle{pack25519} is effectively reducing mod $\p$ and returning a number in base $2^8$.

\begin{Coq}
Lemma Pack25519_mod_25519 :
forall (l:list Z),
Zlength l = 16 ->
Forall (fun x => -2^62 < x < 2^62) l ->
ZofList 8 (Pack25519 l) = (Z16.lst l) mod (2^255-19).
\end{Coq}

\subsection{Using the Verifiable Software Toolchain}
\label{using-VST}
The Verifiable Software Toolchain uses a strongest postcondition strategy.
The user must first write a formal specification of the function he wants to verify in Coq.
This should be as close as possible to the C implementation behavior.
This will simplify the proof and help with stepping throught the CLight version of the software.
With the range of inputes defined, VST steps mechanically through each instruction
and ask the user to verify auxiliary goals such as array bound access, or absence of overflows/underflows.
We call this specification a low level specification. A user will then have an easier
time to prove that his low level specification matches a simpler higher level one.

In order to further speed-up the verification process, it has to be know that to
prove \VSTe{crypto_scalarmult}, a user only need the specification of e.g. \VSTe{M}.
This provide with multiple advantages: the verification by the Coq kernel can be done
in parallel and multiple users can work on proving different functions at the same time.
For the sake of completeness we proved all intermediate functions.

Memory aliasing is the next point a user should pay attention. The way VST
deals with the separation logic is similar to a consumer producer problem.
A simple specification of \texttt{M(o,a,b)} will assume three distinct memory share.
When called with three memory share (\texttt{o, a, b}), the three of them will be consumed.
However assuming this naive specification when \texttt{M(o,a,a)} is called (squaring),
the first two memory shares (\texttt{o, a}) are consumed and VST will expect a third
memory share where the last \texttt{a} is pointing at which does not \textit{exist} anymore.

Examples of such cases are summarized in Fig \ref{tk:MemSame}.
\begin{figure}[h]
  \begin{tikzpicture}
\draw (0,0) rectangle (3,0.3);
\draw[thick] (1,0) -- (1,0.3);
\draw[thick] (2,0) -- (2,0.3);
\node [anchor=east] (shn) at (0,0.15) {sh};
\node [anchor=north] (Moab) at (1.5,-0.5) {\texttt{M(o,a,b)}};
\draw[thick, -> ] (1.25, -0.6) -- (1.25, -0.5) -- (0.5,-0.1);
\draw[thick, -> ] (1.6, -0.6) -- (1.6, -0.5) -- (1.5,-0.1);
\draw[thick, -> ] (1.95, -0.6) -- (1.95, -0.5) -- (2.5,-0.1);

\node [anchor=north west, text width=6cm, align=left] (Sepoab) at (-0.5,-1.8)
  {\VSTess{sh [{ v_o $\!\!\}\!\!]\!\!\!$ <<(lg16)-- o}\\[-0.2em]
  \VSTess{sh [{ v_a $\!\!\}\!\!]\!\!\!$ <<(lg16)-- a}\\[-0.2em]
  \VSTess{sh [{ v_b $\!\!\}\!\!]\!\!\!$ <<(lg16)-- b}};


\begin{scope}[yshift=0 cm,xshift=4.5 cm]
  \draw (0,0) rectangle (2,0.3);
  \draw[thick] (1,0) -- (1,0.3);
  % \draw[thick] (2,0) -- (2,0.3);
  \node [anchor=east] (shm) at (0,0.15) {sh};
  \node [anchor=north] (Mcaa) at (1,-0.5) {\texttt{M(c,c,a)}};
  \draw[thick, -> ] (0.75, -0.6) -- (0.75, -0.5) -- (0.4,-0.1);
  \draw[thick, -> ] (1.1, -0.6) -- (1.1, -0.5) -- (0.6,-0.1);
  \draw[thick, -> ] (1.45, -0.6) -- (1.45, -0.5) -- (1.5,-0.1);

  \node [anchor=north west, text width=6cm, align=left] (Sepoab) at (-0.5,-1.8)
    {\VSTess{sh [{ v_c $\!\!\}\!\!]\!\!\!$ <<(lg16)-- c}\\[-0.2em]
    \VSTess{sh [{ v_a $\!\!\}\!\!]\!\!\!$ <<(lg16)-- a}};
\end{scope}

\draw[dashed] (-0.5,-1.2) -- +(8,0);
\node [anchor=north west] (sep) at (-0.5,-1.2) {In Separation Logic:};

\begin{scope}[yshift=-4.25 cm,xshift=2 cm]
  \draw (0,0) rectangle (2,0.3);
  \draw[thick] (1,0) -- (1,0.3);
  \node [anchor=east] (shm) at (0,0.15) {sh\;};

  \draw (0,0.5) rectangle (1,0.8);
  \node [anchor=east] (shm) at (0,0.65) {Ews};

  \draw[thick] (2,0) -- (2,0.3);
  \node [anchor=north west] (Mcaa) at (0.15,-0.5) {\texttt{M(a,c,\_121665)}};
  \draw[thick, -> ] (0.75, -0.6) -- (0.75, -0.5) -- (0.5,-0.1);
  \draw[thick, -> ] (1.1, -0.6) -- (1.1, -0.5) -- (1.5,-0.1);
  \draw[thick, -> ] (2.0, -0.6) -- (2.45, -0.2) -- (2.45, 0.3) -- (2.2, 0.65) -- (1.5,0.65);

  \node [anchor=north west, text width=6cm, align=left] (Sepoab) at (-0.5,-1.8)
    {\VSTess{sh [{ v_a $\!\!\}\!\!]\!\!\!$ <<(lg16)-- a}\\[-0.2em]
    \VSTess{sh [{ v_c $\!\!\}\!\!]\!\!\!$ <<(lg16)-- c}\\[-0.2em]
    \VSTess{Ews [{ __121665 $\!\!\}\!\!]\!\!\!$ <<(lg16)-- _121665}
    };
\end{scope}
\draw[dashed] (0.5,-5.45) -- +(6,0);
\node [anchor=north west] (sep) at (0.5,-5.45) {In Separation Logic:};
\end{tikzpicture}%

  \caption{Aliasing and Separation Logic}
  \label{tk:MemSame}
\end{figure}

This forces the user to either define multiple specifications for a single function
or specify in his specification which aliasing version is being used.
For our specifications of functions with 3 arguments, named here after \texttt{o, a, b},
we define an additional parameter $k$ with values in
$\{0,1,2,3\}$:
\begin{itemize}
  \item if $k=0$ then \texttt{o} and \texttt{a} are aliased.
  \item if $k=1$ then \texttt{o} and \texttt{b} are aliased.
  \item if $k=2$ then \texttt{a} and \texttt{b} are aliased.
  \item else there is no aliasing.
\end{itemize}
This solution allows us to make cases analysis over possible aliasing.

\subsection{Verifiying \texttt{for} loops}

Final state of \texttt{for} loops are usually computed by simple recursive functions.
However we must define invariants which are true for each iteration.

Assume we want to prove a decreasing loop where indexes go from 3 to 0.
Define a function $g : \N \rightarrow State  \rightarrow State $ which takes as input an integer for the index and a state,
then return a state. It simulate the body of the \texttt{for} loop.
Assume it's recursive call: $f : \N \rightarrow State \rightarrow State $ which iteratively apply $g$ with decreasing index:
\begin{equation*}
  f ( i , s ) =
  \begin{cases}
  s & \text{if } s = 0 \\
  f( i - 1 , g ( i - 1  , s )) & \text{otherwise}
  \end{cases}
\end{equation*}
Then we have :
\begin{align*}
  f(4,s) &= g(0,g(1,g(2,g(3,s))))
  % \\
  % f(3,s) &= g(0,g(1,g(2,s)))
\end{align*}
To prove the correctness of $f(4,s)$, we need to prove that intermediate steps
$g(3,s)$; $g(2,g(3,s))$; $g(1,g(2,g(3,s)))$; $g(0,g(1,g(2,g(3,s))))$ are correct.
Due to the computation order of recursive function, our loop invariant for $i\in\{0;1;2;3;4\}$ cannot use $f(i)$.
To solve this, we define an auxiliary function with an accumulator such that given $i\in\{0;1;2;3;4\}$, it will compute the first $i$ steps of the loop.
We then prove for the complete number of steps, the function with the accumulator and without returns the same result.

We formalized this result in a generic way as follows:
\begin{Coq}
Variable T : Type.
Variable g : nat -> T -> T.

Fixpoint rec_fn (n:nat) (s:T) :=
  match n with
  | 0 => s
  | S n => rec_fn n (g n s)
  end.

Fixpoint rec_fn_rev_acc (n:nat) (m:nat) (s:T) :=
  match n with
  | 0 => s
  | S n => g (m - n - 1) (rec_fn_rev_acc n m s)
  end.

Definition rec_fn_rev (n:nat) (s:T) :=
  rec_fn_rev_acc n n s.

Lemma Tail_Head_equiv :
  forall (n:nat) (s:T),
  rec_fn n s = rec_fn_rev n s.
\end{Coq}
Using this formalization, we prove that the 255 steps of the montgomery ladder in C provide the same computations are the one defined in Algorithm \ref{montgomery-double-add}.
